\documentclass[11pt, a4paper]{article}
\usepackage[a4paper, top=2.5cm, bottom=2.5cm, left=2cm, right=2cm]{geometry}
\usepackage{fontspec}
\usepackage[english, bidi=basic, provide=*]{babel}
\babelprovide[import, onchar=ids fonts]{english}
\babelfont{rm}{TeX Gyre Termes}
\babelfont{sf}{TeX Gyre Heros}
\babelfont{tt}{TeX Gyre Cursor}
\usepackage{enumitem}
\usepackage{amsmath}

\title{Study Notes: Mendel and the Gene Idea (Chapter 14)}
\author{Subject: Biology}
\date{}

\begin{document}
\maketitle

\section*{Core Concepts}
\begin{itemize}
    \item \textbf{Mendel's Particulate Hypothesis:} Parents pass on discrete heritable units (genes) that retain their separate identities in offspring. This replaced the "blending" model.
    \item \textbf{Laws of Inheritance:}
    \begin{itemize}
        \item \textbf{Law of Segregation:} Two alleles for a character separate during gamete formation and end up in different gametes.
        \item \textbf{Law of Independent Assortment:} Each pair of alleles segregates independently of other pairs during gamete formation (applies to genes on different chromosomes).
    \end{itemize}
    \item \textbf{Degrees of Dominance:}
    \begin{itemize}
        \item \textbf{Complete Dominance:} Phenotypes of heterozygote and dominant homozygote are identical.
        \item \textbf{Incomplete Dominance:} F1 hybrids have a phenotype between those of the parents (e.g., pink snapdragons).
        \item \textbf{Codominance:} Two dominant alleles affect the phenotype in separate, distinguishable ways (e.g., AB blood type).
    \end{itemize}
\end{itemize}

\section*{Key Definitions}
\begin{itemize}
    \item \textbf{Allele:} Alternative versions of a gene.
    \item \textbf{Homozygous:} Having two identical alleles for a gene ($PP$ or $pp$).
    \item \textbf{Heterozygous:} Having two different alleles for a gene ($Pp$).
    \item \textbf{Genotype:} Genetic makeup (the alleles present).
    \item \textbf{Phenotype:} Observable physical traits.
    \item \textbf{Testcross:} Breeding an individual of unknown genotype with a homozygous recessive individual to reveal the genotype.
    \item \textbf{Monohybrid vs. Dihybrid:} A cross between individuals heterozygous for one character vs. two characters.
\end{itemize}

\section*{Important Formulas / Relations}
\begin{itemize}
    \item \textbf{Multiplication Rule (AND):} The probability that two or more independent events will occur together is the product of their individual probabilities.
    \item \textbf{Addition Rule (OR):} The probability that any one of two or more exclusive events will occur is the sum of their individual probabilities.
    \item \textbf{Dihybrid Cross Ratio:} $9:3:3:1$ (phenotypic ratio).
    \item \textbf{Gamete Combinations:} $2^n$, where $n$ is the number of heterozygous gene pairs.
\end{itemize}

\section*{Complex Genetic Patterns}
\begin{itemize}
    \item \textbf{Multiple Alleles:} Genes exist in more than two allelic forms (e.g., ABO blood groups with alleles $I^A, I^B, i$).
    \item \textbf{Pleiotropy:} One gene has multiple phenotypic effects (e.g., cystic fibrosis, sickle-cell disease).
    \item \textbf{Epistasis:} A gene at one locus alters the phenotypic expression of a gene at a second locus (e.g., coat color in Labs).
    \item \textbf{Polygenic Inheritance:} An additive effect of two or more genes on a single phenotypic character (e.g., skin color, height).
    \item \textbf{Multifactorial:} Traits that depend on multiple genes combined with environmental influences.
\end{itemize}

\section*{Common Mistakes and Clarifications}
\begin{itemize}
    \item \textbf{Dominance vs. Frequency:} Dominant alleles are not necessarily more common in a population (e.g., polydactyly).
    \item \textbf{Dominance vs. "Strength":} Alleles do not physically interact; dominance is about how genotypes are expressed as phenotypes.
    \item \textbf{Law Applicability:} Independent Assortment only works for genes on different chromosomes or very far apart on the same one.
\end{itemize}

\section*{Human Genetic Disorders}
\begin{itemize}
    \item \textbf{Recessive Disorders:} Require two copies of the allele. \textbf{Carriers} are heterozygous and phenotypically normal. (e.g., Albinism, Cystic Fibrosis, Sickle-cell).
    \item \textbf{Dominant Disorders:} Require only one copy. (e.g., Achondroplasia, Huntington’s Disease).
    \item \textbf{Fetal Testing:} 
    \begin{itemize}
        \item \textbf{Amniocentesis:} Liquid sampled at 15 weeks.
        \item \textbf{CVS:} Tissue sampled from placenta at 10 weeks.
    \end{itemize}
\end{itemize}

\end{document}